\documentclass[11pt,a4paper]{article}

\newcommand{\doctitle}{Flood Polynomial Model (MKM-FPO-001)}
\newcommand{\docsubtitle}{Wolfram Oracle Benchmark of v1.0 and its GBM Approximation}
\newcommand{\docversion}{0.1}
\newcommand{\docdate}{03-November-2026}
\newcommand{\docauthor}{David K Kelly}

% MKM Research Labs — Shared Document Header
% Usage: % MKM Research Labs — Shared Document Header
% Usage: % MKM Research Labs — Shared Document Header
% Usage: \input{../shared/mkm_header}  (before \begin{document})
%
% Required commands before \input:
%   \newcommand{\doctitle}{...}
%   \newcommand{\docsubtitle}{...}
%   \newcommand{\docversion}{1.0}
%   \newcommand{\docdate}{13-February-2026}
%   \newcommand{\docauthor}{David K Kelly}

% ─── Packages ───────────────────────────────────────────────────────────────
\usepackage[utf8]{inputenc}
\usepackage{amsmath,amsfonts,amssymb}
\usepackage{listings}
\usepackage{color,xcolor}
\usepackage{graphicx}
\usepackage{hyperref}
\usepackage[toc,page]{appendix}
\usepackage{tabto}
\usepackage{float}
\usepackage{comment}
\usepackage[en-GB]{datetime2}
\usepackage{setspace}
\usepackage{subcaption}
\usepackage{booktabs}
\usepackage{longtable}
\usepackage{adjustbox}
\usepackage{enumitem}
\usepackage{fancyhdr}
\usepackage{lastpage}

% ─── Page geometry ──────────────────────────────────────────────────────────
\onehalfspacing
\usepackage{geometry}
\geometry{a4paper, total={170mm,247mm}, left=20mm, top=25mm, bottom=25mm}

% ─── Headers and footers ───────────────────────────────────────────────────
\pagestyle{fancy}
\fancyhf{}
\fancyhead[L]{\small\textcolor{gray}{MKM Research Labs}}
\fancyhead[R]{\small\textcolor{gray}{\doctitle{} --- v\docversion}}
\fancyfoot[C]{\small\textcolor{gray}{CONFIDENTIAL}}
\fancyfoot[R]{\small\textcolor{gray}{Page \thepage\ of \pageref{LastPage}}}
\renewcommand{\headrulewidth}{0.4pt}
\renewcommand{\footrulewidth}{0.4pt}

% ─── Colours ────────────────────────────────────────────────────────────────
\definecolor{mkmblue}{HTML}{1565C0}
\definecolor{codegreen}{rgb}{0,0.6,0}
\definecolor{codegray}{rgb}{0.5,0.5,0.5}
\definecolor{codepurple}{rgb}{0.58,0,0.82}
\definecolor{backcolour}{rgb}{0.95,0.95,0.92}

% ─── Hyperref setup ─────────────────────────────────────────────────────────
\hypersetup{
    colorlinks=true,
    linkcolor=mkmblue,
    citecolor=mkmblue,
    urlcolor=mkmblue
}

% ─── Code listings ──────────────────────────────────────────────────────────
\lstdefinestyle{mystyle}{
    backgroundcolor=\color{backcolour},
    commentstyle=\color{codegreen},
    keywordstyle=\color{magenta},
    numberstyle=\tiny\color{codegray},
    stringstyle=\color{codepurple},
    basicstyle=\ttfamily\footnotesize,
    breakatwhitespace=false,
    breaklines=true,
    captionpos=b,
    keepspaces=true,
    numbers=left,
    numbersep=5pt,
    showspaces=false,
    showstringspaces=false,
    showtabs=false,
    tabsize=2
}
\lstset{style=mystyle}

% ─── Section formatting ────────────────────────────────────────────────────
\setcounter{tocdepth}{3}
\setcounter{secnumdepth}{3}

% ─── Shared macros ──────────────────────────────────────────────────────────
\newcommand{\mkmtitlepage}{%
    \begin{titlepage}
    \begin{center}
    \vspace*{1cm}
    \begin{figure}[H]\centering
        \includegraphics[width=0.4\textwidth]{../shared/Logos/MKM_logo.png}
    \end{figure}
    \vspace{1.5cm}
    {\Huge\bfseries\textcolor{mkmblue}{\doctitle}\par}
    \vspace{0.5cm}
    {\Large\textcolor{gray}{\docsubtitle}\par}
    \vspace{2cm}
    {\large MKM Research Labs\par}
    \vspace{0.5cm}
    {\large Version \docversion{} --- \docdate\par}
    \vspace{0.5cm}
    {\large \docauthor\par}
    \vspace{2cm}
    \textcolor{red}{\textbf{CONFIDENTIAL --- PROPRIETARY}}
    \end{center}
    \end{titlepage}
}

\newcommand{\mkmlegalpage}{%
    \clearpage
    \thispagestyle{empty}
    \begin{center}
    \large\textbf{Legal Notice}
    \vspace{1em}

    \textbf{PROPRIETARY AND CONFIDENTIAL}
    \end{center}

    \noindent This document, including all algorithms, models, methodology, formulas,
    calculations, and intellectual concepts contained herein, constitutes the exclusive
    intellectual property and confidential trade secrets of MKM Research Labs (``MKM'').

    \vspace{0.5em}
    \noindent Any usage, reproduction, distribution, modification, or disclosure of this
    document or any portion thereof, without the express prior written authorisation from
    MKM Research Labs is strictly prohibited and constitutes an infringement of intellectual
    property rights.

    \vspace{0.5em}
    \noindent The document is provided on an ``as is'' basis. MKM Research Labs makes no
    representations or warranties of any kind, express or implied, regarding its accuracy,
    reliability, or suitability for any particular purpose.

    \vspace{1em}
    \begin{center}
    All rights reserved. \copyright\ 2021--2026 MKM Research Labs.

    \vspace{0.5em}
    Governed by the laws of the United Kingdom.
    \end{center}
    \clearpage
}

\newcommand{\mkmhistorypage}[1]{%
    \section{Document History}
    \begin{table}[ht]
        \centering
        \begin{tabular}{l|l|l|l}
            \textbf{Date} & \textbf{Version} & \textbf{Description} & \textbf{Author} \\
            \hline
            #1
        \end{tabular}
        \label{tab:revision_history}
    \end{table}
    \clearpage
}

% ─── Bank validation protocol section commands ──────────────────────────────
% Use these to ensure every model doc has the required sections.
% They are just styled \section calls — content is provided by each doc.

\newcommand{\mkmexecsummary}{\section{Executive Summary}}
\newcommand{\mkmpurpose}{\section{Model Purpose and Scope}}
\newcommand{\mkmframework}{\section{Mathematical Framework}}
\newcommand{\mkminputs}{\section{Input Parameters and Calibration}}
\newcommand{\mkmimplementation}{\section{Implementation Details}}
\newcommand{\mkmsensitivity}{\section{Sensitivity Analysis}}
\newcommand{\mkmvalidation}{\section{Validation and Backtesting}}
\newcommand{\mkmtestresults}{\subsection{Automated Test Results}\input{test_results}}
\newcommand{\mkmlimitations}{\section{Model Limitations and Known Weaknesses}}

\newcommand{\mkmgovernance}{\section{Governance Validation Questions}
\noindent The following nine questions are mandated by Chapter~5 of the
\emph{Model Risk Governance for Vendors} handbook and must be explicitly
addressed for every model in the inventory. Response status is tracked in the
Model Governance Dashboard (MRC portal). \par\medskip}

\newcommand{\mkmvqitem}[3]{\subsection{Q#1 --- #2}
\noindent\textit{``#3''}\par\medskip\noindent}

\newcommand{\mkmchangelog}{\section{Change History}}
  (before \begin{document})
%
% Required commands before \input:
%   \newcommand{\doctitle}{...}
%   \newcommand{\docsubtitle}{...}
%   \newcommand{\docversion}{1.0}
%   \newcommand{\docdate}{13-February-2026}
%   \newcommand{\docauthor}{David K Kelly}

% ─── Packages ───────────────────────────────────────────────────────────────
\usepackage[utf8]{inputenc}
\usepackage{amsmath,amsfonts,amssymb}
\usepackage{listings}
\usepackage{color,xcolor}
\usepackage{graphicx}
\usepackage{hyperref}
\usepackage[toc,page]{appendix}
\usepackage{tabto}
\usepackage{float}
\usepackage{comment}
\usepackage[en-GB]{datetime2}
\usepackage{setspace}
\usepackage{subcaption}
\usepackage{booktabs}
\usepackage{longtable}
\usepackage{adjustbox}
\usepackage{enumitem}
\usepackage{fancyhdr}
\usepackage{lastpage}

% ─── Page geometry ──────────────────────────────────────────────────────────
\onehalfspacing
\usepackage{geometry}
\geometry{a4paper, total={170mm,247mm}, left=20mm, top=25mm, bottom=25mm}

% ─── Headers and footers ───────────────────────────────────────────────────
\pagestyle{fancy}
\fancyhf{}
\fancyhead[L]{\small\textcolor{gray}{MKM Research Labs}}
\fancyhead[R]{\small\textcolor{gray}{\doctitle{} --- v\docversion}}
\fancyfoot[C]{\small\textcolor{gray}{CONFIDENTIAL}}
\fancyfoot[R]{\small\textcolor{gray}{Page \thepage\ of \pageref{LastPage}}}
\renewcommand{\headrulewidth}{0.4pt}
\renewcommand{\footrulewidth}{0.4pt}

% ─── Colours ────────────────────────────────────────────────────────────────
\definecolor{mkmblue}{HTML}{1565C0}
\definecolor{codegreen}{rgb}{0,0.6,0}
\definecolor{codegray}{rgb}{0.5,0.5,0.5}
\definecolor{codepurple}{rgb}{0.58,0,0.82}
\definecolor{backcolour}{rgb}{0.95,0.95,0.92}

% ─── Hyperref setup ─────────────────────────────────────────────────────────
\hypersetup{
    colorlinks=true,
    linkcolor=mkmblue,
    citecolor=mkmblue,
    urlcolor=mkmblue
}

% ─── Code listings ──────────────────────────────────────────────────────────
\lstdefinestyle{mystyle}{
    backgroundcolor=\color{backcolour},
    commentstyle=\color{codegreen},
    keywordstyle=\color{magenta},
    numberstyle=\tiny\color{codegray},
    stringstyle=\color{codepurple},
    basicstyle=\ttfamily\footnotesize,
    breakatwhitespace=false,
    breaklines=true,
    captionpos=b,
    keepspaces=true,
    numbers=left,
    numbersep=5pt,
    showspaces=false,
    showstringspaces=false,
    showtabs=false,
    tabsize=2
}
\lstset{style=mystyle}

% ─── Section formatting ────────────────────────────────────────────────────
\setcounter{tocdepth}{3}
\setcounter{secnumdepth}{3}

% ─── Shared macros ──────────────────────────────────────────────────────────
\newcommand{\mkmtitlepage}{%
    \begin{titlepage}
    \begin{center}
    \vspace*{1cm}
    \begin{figure}[H]\centering
        \includegraphics[width=0.4\textwidth]{../shared/Logos/MKM_logo.png}
    \end{figure}
    \vspace{1.5cm}
    {\Huge\bfseries\textcolor{mkmblue}{\doctitle}\par}
    \vspace{0.5cm}
    {\Large\textcolor{gray}{\docsubtitle}\par}
    \vspace{2cm}
    {\large MKM Research Labs\par}
    \vspace{0.5cm}
    {\large Version \docversion{} --- \docdate\par}
    \vspace{0.5cm}
    {\large \docauthor\par}
    \vspace{2cm}
    \textcolor{red}{\textbf{CONFIDENTIAL --- PROPRIETARY}}
    \end{center}
    \end{titlepage}
}

\newcommand{\mkmlegalpage}{%
    \clearpage
    \thispagestyle{empty}
    \begin{center}
    \large\textbf{Legal Notice}
    \vspace{1em}

    \textbf{PROPRIETARY AND CONFIDENTIAL}
    \end{center}

    \noindent This document, including all algorithms, models, methodology, formulas,
    calculations, and intellectual concepts contained herein, constitutes the exclusive
    intellectual property and confidential trade secrets of MKM Research Labs (``MKM'').

    \vspace{0.5em}
    \noindent Any usage, reproduction, distribution, modification, or disclosure of this
    document or any portion thereof, without the express prior written authorisation from
    MKM Research Labs is strictly prohibited and constitutes an infringement of intellectual
    property rights.

    \vspace{0.5em}
    \noindent The document is provided on an ``as is'' basis. MKM Research Labs makes no
    representations or warranties of any kind, express or implied, regarding its accuracy,
    reliability, or suitability for any particular purpose.

    \vspace{1em}
    \begin{center}
    All rights reserved. \copyright\ 2021--2026 MKM Research Labs.

    \vspace{0.5em}
    Governed by the laws of the United Kingdom.
    \end{center}
    \clearpage
}

\newcommand{\mkmhistorypage}[1]{%
    \section{Document History}
    \begin{table}[ht]
        \centering
        \begin{tabular}{l|l|l|l}
            \textbf{Date} & \textbf{Version} & \textbf{Description} & \textbf{Author} \\
            \hline
            #1
        \end{tabular}
        \label{tab:revision_history}
    \end{table}
    \clearpage
}

% ─── Bank validation protocol section commands ──────────────────────────────
% Use these to ensure every model doc has the required sections.
% They are just styled \section calls — content is provided by each doc.

\newcommand{\mkmexecsummary}{\section{Executive Summary}}
\newcommand{\mkmpurpose}{\section{Model Purpose and Scope}}
\newcommand{\mkmframework}{\section{Mathematical Framework}}
\newcommand{\mkminputs}{\section{Input Parameters and Calibration}}
\newcommand{\mkmimplementation}{\section{Implementation Details}}
\newcommand{\mkmsensitivity}{\section{Sensitivity Analysis}}
\newcommand{\mkmvalidation}{\section{Validation and Backtesting}}
\newcommand{\mkmtestresults}{\subsection{Automated Test Results}\input{test_results}}
\newcommand{\mkmlimitations}{\section{Model Limitations and Known Weaknesses}}

\newcommand{\mkmgovernance}{\section{Governance Validation Questions}
\noindent The following nine questions are mandated by Chapter~5 of the
\emph{Model Risk Governance for Vendors} handbook and must be explicitly
addressed for every model in the inventory. Response status is tracked in the
Model Governance Dashboard (MRC portal). \par\medskip}

\newcommand{\mkmvqitem}[3]{\subsection{Q#1 --- #2}
\noindent\textit{``#3''}\par\medskip\noindent}

\newcommand{\mkmchangelog}{\section{Change History}}
  (before \begin{document})
%
% Required commands before \input:
%   \newcommand{\doctitle}{...}
%   \newcommand{\docsubtitle}{...}
%   \newcommand{\docversion}{1.0}
%   \newcommand{\docdate}{13-February-2026}
%   \newcommand{\docauthor}{David K Kelly}

% ─── Packages ───────────────────────────────────────────────────────────────
\usepackage[utf8]{inputenc}
\usepackage{amsmath,amsfonts,amssymb}
\usepackage{listings}
\usepackage{color,xcolor}
\usepackage{graphicx}
\usepackage{hyperref}
\usepackage[toc,page]{appendix}
\usepackage{tabto}
\usepackage{float}
\usepackage{comment}
\usepackage[en-GB]{datetime2}
\usepackage{setspace}
\usepackage{subcaption}
\usepackage{booktabs}
\usepackage{longtable}
\usepackage{adjustbox}
\usepackage{enumitem}
\usepackage{fancyhdr}
\usepackage{lastpage}

% ─── Page geometry ──────────────────────────────────────────────────────────
\onehalfspacing
\usepackage{geometry}
\geometry{a4paper, total={170mm,247mm}, left=20mm, top=25mm, bottom=25mm}

% ─── Headers and footers ───────────────────────────────────────────────────
\pagestyle{fancy}
\fancyhf{}
\fancyhead[L]{\small\textcolor{gray}{MKM Research Labs}}
\fancyhead[R]{\small\textcolor{gray}{\doctitle{} --- v\docversion}}
\fancyfoot[C]{\small\textcolor{gray}{CONFIDENTIAL}}
\fancyfoot[R]{\small\textcolor{gray}{Page \thepage\ of \pageref{LastPage}}}
\renewcommand{\headrulewidth}{0.4pt}
\renewcommand{\footrulewidth}{0.4pt}

% ─── Colours ────────────────────────────────────────────────────────────────
\definecolor{mkmblue}{HTML}{1565C0}
\definecolor{codegreen}{rgb}{0,0.6,0}
\definecolor{codegray}{rgb}{0.5,0.5,0.5}
\definecolor{codepurple}{rgb}{0.58,0,0.82}
\definecolor{backcolour}{rgb}{0.95,0.95,0.92}

% ─── Hyperref setup ─────────────────────────────────────────────────────────
\hypersetup{
    colorlinks=true,
    linkcolor=mkmblue,
    citecolor=mkmblue,
    urlcolor=mkmblue
}

% ─── Code listings ──────────────────────────────────────────────────────────
\lstdefinestyle{mystyle}{
    backgroundcolor=\color{backcolour},
    commentstyle=\color{codegreen},
    keywordstyle=\color{magenta},
    numberstyle=\tiny\color{codegray},
    stringstyle=\color{codepurple},
    basicstyle=\ttfamily\footnotesize,
    breakatwhitespace=false,
    breaklines=true,
    captionpos=b,
    keepspaces=true,
    numbers=left,
    numbersep=5pt,
    showspaces=false,
    showstringspaces=false,
    showtabs=false,
    tabsize=2
}
\lstset{style=mystyle}

% ─── Section formatting ────────────────────────────────────────────────────
\setcounter{tocdepth}{3}
\setcounter{secnumdepth}{3}

% ─── Shared macros ──────────────────────────────────────────────────────────
\newcommand{\mkmtitlepage}{%
    \begin{titlepage}
    \begin{center}
    \vspace*{1cm}
    \begin{figure}[H]\centering
        \includegraphics[width=0.4\textwidth]{../shared/Logos/MKM_logo.png}
    \end{figure}
    \vspace{1.5cm}
    {\Huge\bfseries\textcolor{mkmblue}{\doctitle}\par}
    \vspace{0.5cm}
    {\Large\textcolor{gray}{\docsubtitle}\par}
    \vspace{2cm}
    {\large MKM Research Labs\par}
    \vspace{0.5cm}
    {\large Version \docversion{} --- \docdate\par}
    \vspace{0.5cm}
    {\large \docauthor\par}
    \vspace{2cm}
    \textcolor{red}{\textbf{CONFIDENTIAL --- PROPRIETARY}}
    \end{center}
    \end{titlepage}
}

\newcommand{\mkmlegalpage}{%
    \clearpage
    \thispagestyle{empty}
    \begin{center}
    \large\textbf{Legal Notice}
    \vspace{1em}

    \textbf{PROPRIETARY AND CONFIDENTIAL}
    \end{center}

    \noindent This document, including all algorithms, models, methodology, formulas,
    calculations, and intellectual concepts contained herein, constitutes the exclusive
    intellectual property and confidential trade secrets of MKM Research Labs (``MKM'').

    \vspace{0.5em}
    \noindent Any usage, reproduction, distribution, modification, or disclosure of this
    document or any portion thereof, without the express prior written authorisation from
    MKM Research Labs is strictly prohibited and constitutes an infringement of intellectual
    property rights.

    \vspace{0.5em}
    \noindent The document is provided on an ``as is'' basis. MKM Research Labs makes no
    representations or warranties of any kind, express or implied, regarding its accuracy,
    reliability, or suitability for any particular purpose.

    \vspace{1em}
    \begin{center}
    All rights reserved. \copyright\ 2021--2026 MKM Research Labs.

    \vspace{0.5em}
    Governed by the laws of the United Kingdom.
    \end{center}
    \clearpage
}

\newcommand{\mkmhistorypage}[1]{%
    \section{Document History}
    \begin{table}[ht]
        \centering
        \begin{tabular}{l|l|l|l}
            \textbf{Date} & \textbf{Version} & \textbf{Description} & \textbf{Author} \\
            \hline
            #1
        \end{tabular}
        \label{tab:revision_history}
    \end{table}
    \clearpage
}

% ─── Bank validation protocol section commands ──────────────────────────────
% Use these to ensure every model doc has the required sections.
% They are just styled \section calls — content is provided by each doc.

\newcommand{\mkmexecsummary}{\section{Executive Summary}}
\newcommand{\mkmpurpose}{\section{Model Purpose and Scope}}
\newcommand{\mkmframework}{\section{Mathematical Framework}}
\newcommand{\mkminputs}{\section{Input Parameters and Calibration}}
\newcommand{\mkmimplementation}{\section{Implementation Details}}
\newcommand{\mkmsensitivity}{\section{Sensitivity Analysis}}
\newcommand{\mkmvalidation}{\section{Validation and Backtesting}}
\newcommand{\mkmtestresults}{\subsection{Automated Test Results}\input{test_results}}
\newcommand{\mkmlimitations}{\section{Model Limitations and Known Weaknesses}}

\newcommand{\mkmgovernance}{\section{Governance Validation Questions}
\noindent The following nine questions are mandated by Chapter~5 of the
\emph{Model Risk Governance for Vendors} handbook and must be explicitly
addressed for every model in the inventory. Response status is tracked in the
Model Governance Dashboard (MRC portal). \par\medskip}

\newcommand{\mkmvqitem}[3]{\subsection{Q#1 --- #2}
\noindent\textit{``#3''}\par\medskip\noindent}

\newcommand{\mkmchangelog}{\section{Change History}}

% T1 so the underscores in code identifiers print as underscores.
\usepackage[T1]{fontenc}
\usepackage{lmodern}
% Generated by tools.wolfram_oracle -- do not edit.
\newcommand{\bmGenerated}{2026-11-03}
\newcommand{\bmCodeCommit}{\texttt{5ef357e130c7}}
\newcommand{\bmCodeDirty}{}
\newcommand{\bmNumerics}{numpy 1.26.4, scipy 1.15.1, mpmath 1.4.1}
\newcommand{\bmWolfram}{15.0.0 for Mac OS X ARM (64-bit) (May 26, 2026)}
\newcommand{\bmWolframVia}{ via live kernel session (wolframclient)}
\newcommand{\bmWolframConfirms}{Wolfram, from the exact coefficients and knots, agrees with the reference to $3.1\times10^{-31}$.}
\newcommand{\bmChecks}{10}
\newcommand{\bmOutcome}{all pass}
\newcommand{\bmClaims}{11}
\newcommand{\bmClaimsMet}{11}
\newcommand{\bmGauges}{11}
\newcommand{\bmNearMae}{0.043 to 0.069}
\newcommand{\bmNearOver}{11 to 22}
\newcommand{\bmNearBias}{-0.033 to +0.028}
\newcommand{\bmBelowMae}{0.0066 to 0.0217}
\newcommand{\bmGridMatched}{88}
\newcommand{\bmGridCells}{88}
\newcommand{\bmHeld}{23}
\newcommand{\bmDrop}{0.0e+00}
\newcommand{\bmEnvWorst}{0.05}
\newcommand{\bmMinFree}{0.0057}
\newcommand{\bmMinEnv}{0.0105}

% The shared history table sets every column to its natural width; wrap the description.
\renewcommand{\mkmhistorypage}[1]{%
    \section{Document History}
    \begin{table}[ht]
        \centering
        \begin{tabular}{l|l|p{9cm}|l}
            \textbf{Date} & \textbf{Version} & \textbf{Description} & \textbf{Author} \\
            \hline
            #1
        \end{tabular}
    \end{table}
}

\begin{document}
% Long code paths and identifiers cannot hyphenate; let the spacing stretch instead.
\sloppy

\mkmtitlepage
\mkmlegalpage

\tableofcontents
\clearpage

% =========================================================================
\mkmexecsummary
% =========================================================================

This document benchmarks the Flood Polynomial Model (MKM-FPO-001), v1.0, a closed-form
approximation of the per-gauge GBM flood classifiers (MKM-FC-001) that the
single-gauge stress scenario uses to price PRS trades hour by hour before a breach. It
asks whether the code computes what its model document (\texttt{flood\_poly.pdf}, v1.0)
says, whether the document's claims hold, and how well the polynomial approximates the
classifiers it stands in for.

\paragraph{The code computes its documented formula.} \texttt{p\_flood\_poly} agrees
with a 50-digit implementation of Eqs.~1--4 and the printed coefficients, and so do
\texttt{p\_flood\_series}, the copy in \texttt{stress.sensitivity} and its slope. The
port stress screen's JavaScript holds the same coefficients
(Table~\ref{tab:bm-checks}). \bmWolframConfirms

\paragraph{But the document's evidence does not hold} (Table~\ref{tab:bm-claims}).
\begin{itemize}[leftmargin=*,nosep]
\item \textbf{Its tables have no source.} The code reproduces \bmTableOne{} of the 54
  cells of Table~1 and none of Table~2. Nor does the GBM reproduce them. No generator
  for either table ever existed.
\item \textbf{At the severe level, $P$ does not rise from 0.42 to 0.97.} It is 0.00 at
  hour~0 and 0.29 at hour~167, after peaking near 0.97 around hour~96.
\item \textbf{The fit is not as stated.} On every real storm state of the fitted gauge,
  $R^2$ is \bmFitRsq, not 0.9395. MAE is \bmFitMae, better than the stated 0.0456. The
  claimed 0.05 accuracy is not met near severe.
\end{itemize}

\paragraph{The approximation, measured.} The comparison uses the \bmGauges{} gauges'
classifiers on the real storm states they are trained on (Table~\ref{tab:bm-gauges}).
\begin{itemize}[leftmargin=*,nosep]
\item Below severe, the polynomial is within \bmBelowMae{} of the GBM on average.
\item Near severe (75--100\%), where stress P\&L is decided, the mean error is
  \bmNearMae, and \bmNearOver\% of states are off by more than 0.10.
\item The bias near severe ranges from \bmNearBias. It is right on average but
  imprecise state by state.
\item Early in a storm the surface is not monotone in the level: at hour~24, $P$ passes
  $\tfrac12$ twice below severe and falls as the water rises between the two levels.
\end{itemize}
Two things production does are not in the document. The stress route latches $P=1$ once
the level reaches severe. And \texttt{p\_flood\_simple}, which sets portfolio stress P\&L,
is the ratio $\min(1,w/s)$, not a probability.

\paragraph{Evidence.} \bmChecks{} checks, outcome: \bmOutcome{}
(Table~\ref{tab:bm-checks}).
\begin{itemize}[leftmargin=*,nosep]
\item Wolfram: \bmWolfram{}\bmWolframVia.
\item Code under test, at commit \bmCodeCommit{}\bmCodeDirty.
\item Numerics: \bmNumerics.
\item Generated \bmGenerated{} by \texttt{python -m tools.wolfram\_oracle fpo001}.
\end{itemize}

% =========================================================================
\mkmpurpose
% =========================================================================

\subsection{What is benchmarked}
\begin{itemize}[leftmargin=*,nosep]
\item \texttt{src/models/stress/flood\_poly.py}: \texttt{p\_flood\_poly},
  \texttt{p\_flood\_series} and \texttt{p\_flood\_simple}.
\item The copy in \texttt{src/models/stress/sensitivity.py}.
\item The coefficients hard-coded in the port stress screen.
\item The polynomial against the production GBM classifier of every gauge that has one.
\end{itemize}

\subsection{What is not}
\begin{itemize}[leftmargin=*,nosep]
\item The classifiers themselves (MKM-FC-001).
\item The stress scenarios' hydrographs.
\item The P\&L built on $P$.
\end{itemize}
This is a benchmark of the implementation and of the approximation (validation question
Q8), not independent validation (Q9).

% =========================================================================
\mkmframework
% =========================================================================

The document's model is
\begin{equation}
P=\sigma(a h+b t+c h t+d h^2+e t^2+f),\quad h=\ln\max(w/s,10^{-8}),\quad
t=\ln\frac{\text{hour}+1}{168},
\end{equation}
with the exponent clamped to $\pm30$. The reference evaluates it at 50 digits from the
printed coefficients.

\paragraph{What the polynomial can know.} Each gauge's GBM also sees the rate of rise
$\Delta h$ and its change. A function of $(h,t)$ alone can at best reproduce the GBM's
average over the storm states that share a level and an hour. The approximation is
therefore judged on those real states, as production meets them, not on a grid at
$\Delta h=0$. A grid includes states storms do not produce, and states storms do produce
are left out of it.

% =========================================================================
\section{Benchmark Method}
% =========================================================================

\paragraph{The code.} The code is checked against the reference over three severe
thresholds, sixteen level ratios from zero to 1.5, and thirteen hours.

\paragraph{Wolfram.} From exact coefficients, Wolfram evaluates the clamped logistic at
every cell of Table~1, the levels where $P=\tfrac12$ (by \texttt{NSolve} over the reals
below severe), and $\partial P/\partial h$ symbolically.

\paragraph{The GBM comparison.}
\begin{itemize}[leftmargin=*,nosep]
\item For every gauge with a classifier, 2{,}000 storms are drawn once (seed 7) from the
  catalogue. Each storm's states are rebuilt through production's training-vector
  builder.
\item The polynomial and the GBM are compared on the states below severe.
\item For the fitted gauge, the fit metrics are recomputed on every state of the full
  catalogue. The cubic of \S9.1 is fitted to the same states.
\end{itemize}

% =========================================================================
\mkmvalidation
% =========================================================================

% Generated by tools.wolfram_oracle -- do not edit.
\begin{table}[H]\centering\small
\caption{Benchmark checks}\label{tab:bm-checks}
\begin{adjustbox}{max width=\textwidth}
\begin{tabular}{p{92mm}p{16mm}rl}
\toprule
Check & Result & Limit & Outcome \\
\midrule
Eq.~3: \texttt{p\_flood\_poly} against the reference, worst relative error & $3.8\times10^{-14}$ & $1.0\times10^{-12}$ & pass \\
A threshold at or below zero gives zero & yes &  & pass \\
\texttt{p\_flood\_series} is \texttt{p\_flood\_poly} hour by hour & 0 & 0 & pass \\
\texttt{p\_flood\_simple} is $\min(1, w/s)$ & $7.3\times10^{-17}$ & $1.0\times10^{-12}$ & pass \\
The model as \texttt{stress.sensitivity} evaluates it & $3.8\times10^{-14}$ & $1.0\times10^{-12}$ & pass \\
Its closed-form $\partial P/\partial h$, held hours included & $5.1\times10^{-13}$ & $1.0\times10^{-12}$ & pass \\
The server sends the screen the document's coefficients and envelope & yes &  & pass \\
The screen carries no coefficients of its own & yes &  & pass \\
The screen's JavaScript, run under node, against the reference & $6.8\times10^{-14}$ & $1.0\times10^{-12}$ & pass \\
Wolfram agrees with the reference (table cells, slopes, minima at their points) & $3.1\times10^{-31}$ & $1.0\times10^{-25}$ & pass \\
\bottomrule
\end{tabular}
\end{adjustbox}
\end{table}


\subsection{The gauges}
% Generated by tools.wolfram_oracle -- do not edit.
\begin{table}[H]\centering\small
\caption{The polynomial against each gauge's GBM on real storm states, near severe ($w/s\ge0.75$) and below it}\label{tab:bm-gauges}
\begin{adjustbox}{max width=\textwidth}
\begin{tabular}{lrrrrrrr}
\toprule
Gauge & States & GBM & Poly & Mean $|d|$ & 95\% & $>0.1$ & All below \\
\midrule
GAUGE-b2dbe58b & 14{,}592 & 0.326 & 0.367 & 0.071 & 0.249 & 26\% & 0.0150 \\
GAUGE-dc5e88fb & 9{,}894 & 0.296 & 0.305 & 0.052 & 0.186 & 15\% & 0.0098 \\
GAUGE-1aaf8b08 & 18{,}815 & 0.413 & 0.436 & 0.061 & 0.210 & 21\% & 0.0228 \\
GAUGE-d208bf3d & 14{,}990 & 0.445 & 0.424 & 0.061 & 0.207 & 19\% & 0.0225 \\
GAUGE-c0fbc443 & 10{,}256 & 0.259 & 0.295 & 0.061 & 0.233 & 21\% & 0.0096 \\
GAUGE-a6be4ee9 & 11{,}406 & 0.300 & 0.328 & 0.063 & 0.226 & 22\% & 0.0111 \\
GAUGE-ef53398c & 8{,}379 & 0.281 & 0.291 & 0.052 & 0.209 & 17\% & 0.0078 \\
GAUGE-085522ee & 16{,}268 & 0.416 & 0.421 & 0.060 & 0.203 & 19\% & 0.0186 \\
GAUGE-9fffb377 & 12{,}730 & 0.287 & 0.335 & 0.073 & 0.264 & 27\% & 0.0119 \\
GAUGE-f8629444 & 14{,}105 & 0.406 & 0.395 & 0.059 & 0.197 & 18\% & 0.0182 \\
GAUGE-9e5f3ef9 & 17{,}404 & 0.374 & 0.407 & 0.068 & 0.244 & 24\% & 0.0199 \\
\bottomrule
\end{tabular}
\end{adjustbox}
\end{table}


\subsection{The model document's claims}
% Generated by tools.wolfram_oracle -- do not edit.
\begin{table}[H]\centering\small
\caption{The model document's claims, tested}\label{tab:bm-claims}
\begin{adjustbox}{max width=\textwidth}
\begin{tabular}{p{18mm}p{46mm}p{72mm}l}
\toprule
Claim & Stated & Benchmark finds & Outcome \\
\midrule
\S8, Table~1 & $P$ by level and hour, $s=6.58$\,m & 8 of 54 cells reproduced by the code & \textbf{not met} \\
\S8, Table~2 & GBM and polynomial at six points & polynomial column 0 of 6; GBM column 0 of 6 (rate of rise zero) & \textbf{not met} \\
\S8 & at the severe level, 0.42 at hour~0 rising to 0.97 at hour~167 & 0.00 at hour~0, 0.29 at hour~167 & \textbf{not met} \\
\S8 & $|\text{poly}-\text{GBM}|<0.05$ across the surface & near severe, 95th percentile up to 0.26; above 0.05 in up to 39\% of states & \textbf{not met} \\
\S4.2 & $R^2$ 0.9395, MAE 0.0456, RMSE 0.0986 against the GBM & 0.9083{,} 0.0103{,} 0.0444 on all 1{,}680{,}000 real states & \textbf{not met} \\
\S9.1 & a cubic kernel reaches $R^2\approx0.97$ & 0.9460 on the same states & \textbf{not met} \\
\S4.2 & mildly conservative: over-predicts near severe & bias near severe -0.020 to +0.048 across gauges & \textbf{not met} \\
\S3.2 & $t\approx-5.1$ at hour~0 & -5.1240 & met \\
--- & rising water raises $P$ & at hour~24, $P$ passes $\tfrac12$ at 52\% and 74\% of severe: it falls as the water rises between them & finding \\
\S5 & the scenario uses $P$ as given & the stress route latches $P=1$ once the level reaches severe; undocumented & finding \\
\S5.2 & \texttt{p\_flood\_simple} for portfolio stress & $\min(1,w/s)$ is a ratio, not a probability, and it sets portfolio stress P\&L & finding \\
\bottomrule
\end{tabular}
\end{adjustbox}
\end{table}


% =========================================================================
\mkmlimitations
% =========================================================================

\begin{enumerate}[leftmargin=*]
\item \textbf{Benchmark, not validation.} The GBMs are the reference here; whether they
  are right is MKM-FC-001's question.
\item \textbf{Remediation.} The coefficients are near the best this form can do on real
  states: refitting the same six terms does not reduce the error near severe. The model
  is refitted as v2.0, pooled across gauges, with a richer kernel constrained to rise
  with the level and $P=1$ at or above severe. Its document is reissued with tables
  generated from the code, and it is benchmarked as v0.2 of this document.
\end{enumerate}

\mkmhistorypage{03-Nov-2026 & 0.1 & First issue: tests v1.0 & David K Kelly \\}

\end{document}
